\section{Bilinear distribution}\label{ti:section}

We retain the groups, ordinary prime bands, and weights of the preceding
section. In particular, every ordinary slot contains a prime, there are
$r_0$ labelled slots in each $Q_j$, and
\[
 s_j=s'2^{-j}L,\qquad 20T\le s_{j_x}<40T,\qquad
 \frac{(r_0-1)c_1}{2}\ge c_2+1.
\]
The smooth allocation and change of variables below follow the Type II
method of \cite{SmoothShifted}. We give the argument for the present
all-prime weight and its parity twist, using
Theorem~\ref{sh:correlation} for the resulting correlations.

\begin{theorem}[Type II distribution]\label{ti:distribution}
Fix $D_*,b_*,C>0$. Let $U,V$ be dyadic scales such that
\[
 UV\asymp x,\qquad x^{b_*}\le U,V\le x^{1-b_*},
\]
and let $K\subset[U,2U)$ and $J\subset[V,2V)$ be intervals. Suppose
$|\alpha_m|,|\beta_n|\le L^C$ and $\alpha_m=0$ unless
$P^-(m)>W$. There is a fixed $B$, depending only on these fixed data
and on the fixed weight and cutoff data, for which the following holds.
Assume that the sequence $\alpha_m\1_{m\in K}$ satisfies
\eqref{sh:discrepancy}; explicitly, assume
\begin{equation}\label{ti:actual-discrepancy}
 \left|\sum_{m\in I}\alpha_m\1_{m\in K}
                    \chi(m)m^{-1+it}\right|\le L^{-B}
 \quad\left(
 \begin{array}{c}
 I\subset[1,x^2]\text{ an interval},\\
 \operatorname{mod}(\chi)\le L^B,\quad |t|\le L^B
 \end{array}\right).
\end{equation}
Then
\begin{equation}\label{ti:bound}
 \sum_{\substack{m\in K,\ n\in J,\ u\ge1\\mn=2u+1}}
       \alpha_m\beta_n A(u)\Psi(u/x)
       \ll xL^{-D_*}.
\end{equation}
The estimate is uniform in the intervals and coefficient sequences
satisfying these assumptions. No discrepancy or smoothness assumption is
imposed on $\beta$.
\end{theorem}

\begin{proof}
Extend the actually restricted sequences by zero, and write them again
as $\alpha$ and $\beta$. Thus $\alpha$ in every subsequent convolution
is the sequence appearing in \eqref{ti:actual-discrepancy}. By
$A=A_0(1-\mathcal E)/2$, it suffices to prove the assertion with
$A_0(u)\varepsilon(u)$, for each of the two choices
\[
                    \varepsilon=1\quad\hbox{or}\quad
                    \varepsilon=\mathcal E.
\]
The choice will be the same at the two endpoints of each correlation.

\paragraph{Removing a large group-prime part.}
Choose a fixed band index $j_{**}$ so large that
\begin{equation}\label{ti:designated-range}
 \tau=c_1s'2^{-j_{**}}>0,\qquad
 \eta=c_2s'2^{-j_{**}}<\min(b_*/4,1/4).
\end{equation}
For sufficiently large $x$, this band is present. Designate one of its
labelled slots, whose prime therefore lies in $[x^\tau,x^\eta]$.
This merely names an existing slot in the weight.

Put $u_P=u/u_*$. We first discard $u_P>x^{b_*/4}$. For $u\asymp x$,
the part of $2u+1$ supported on primes exceeding $W$ has at most
$O(L/\log W)=O(\sqrt L)$ prime factors with multiplicity. Its number
of divisors is consequently at most $\exp(O(\sqrt L))$, since
$a+1\le2^a$ for $a\ge1$. This bounds the possible rough divisors $m$.
The tail bound in Lemma~\ref{wt:tail}, with $d=2/5$, gives
\begin{align}
 &\sum_{\substack{mn=2u+1,\ u\asymp x\\u_P>x^{b_*/4}}}
          |\alpha_m\beta_n|A_0(u)|\Psi(u/x)|\notag\\
 &\qquad\ll
 xL^{O(1)}\exp(O(\sqrt L))
   \left(\sum_v\frac{a(v)}v\right)
   \sum_{\substack{r>x^{b_*/4}\\r_*=1}}\frac{W_1(r)}r
 \ll xL^{O(1)}
       \exp\left(-\frac{b_*}{4}L^{1-d}+O(\sqrt L)\right).
 \label{ti:large-group-error}
\end{align}
To obtain the first inequality, write $u=vr$, use $vr\asymp x$ to
bound $1\ll x/(vr)$, and then drop the product restriction.
The harmonic ordinary-slot mass is $H_Q=L^{O(1)}$. Since $1-d>1/2$,
\eqref{ti:large-group-error} is $O_H(xL^{-H})$ for every fixed $H$.
The same conclusion will hold whenever we reinstate these discarded
tuples with a multiplier of absolute value at most one.

\paragraph{A smooth allocation with bounded residual.}
Set
\begin{equation}\label{ti:E}
                            E=UL^{-K_0},
\end{equation}
where the fixed positive constant $K_0$ will be chosen below. Put the
entire group-prime part and the designated prime into $h$. Process all
remaining ordinary slots by increasing band index, in a fixed order
within each band, assigning each slot either to $e$ or to $h$.
Write these slots as $p_1,\ldots,p_N$, with band indices $j(i)$;
here $N=r_0(j_x+1)-1=O(T)$.

Fix a smooth function $\gamma:\mathbb R\to[0,1]$ equal to zero on
$(-\infty,0]$ and to one on $[1,\infty)$. At slot $i$, with residual
target $R_i$, assign the prime to $e$ with weight
\[
        \gamma\left(\frac{R_i-\log p_i}{s_{j(i)}}\right)
\]
and to $h$ with the complementary weight. Start with
$R_1=\log E$, and subtract $\log p_i$ from the residual precisely
when the slot is assigned to $e$.

For a tuple surviving the preceding deletion, the product withheld
from allocation is at most $x^{b_*/2}$. Since $u\asymp x$ and
$U\le x^{1-b_*}$, its available logarithms satisfy, for large $x$,
\[
 0<\log E\le\sum_{i=1}^N\log p_i.
\]
Every branch of positive weight has the invariant
\begin{equation}\label{ti:allocation-invariant}
 0\le R_i\le\sum_{t=i}^N\log p_t+(c_2+1)s_{j_x}
                    \qquad(1\le i\le N+1).
\end{equation}
Indeed, taking a slot requires $R_i>\log p_i$, so preserves both
nonnegativity and the upper bound after subtracting that logarithm.
Skipping a slot in band $j$ requires
\[
                  R_i<\log p_i+s_j\le(c_2+1)s_j.
\]
If $j<j_x$, the next band has at least $r_0-1$ available slots and
therefore total logarithm at least
$(r_0-1)c_1s_j/2\ge(c_2+1)s_j$. These slots are all still
unprocessed. If $j=j_x$, the error term in
\eqref{ti:allocation-invariant} supplies the bound directly.
This proves the invariant by induction, including the terminal step.

The geometric sum of all unprocessed band scales is
$O(s_{j(i)})$. Thus \eqref{ti:allocation-invariant} also places every
argument $(R_i-\log p_i)/s_{j(i)}$ in one fixed compact interval,
independently of $K_0$. At termination,
\[
       0\le\log(E/e)=R_{N+1}\le(c_2+1)s_{j_x}.
\]
Choose once and for all $C_s>40(c_2+1)$. Every surviving branch has
\begin{equation}\label{ti:e-range}
                         EL^{-C_s}\le e\le E.
\end{equation}
This choice of $C_s$ is independent of $K_0$.

Choose a fixed smooth $\rho:\mathbb R\to[0,1]$ of compact support
which equals one throughout the possible argument interval, and put
\[
              f_1=\rho\gamma,\qquad f_0=\rho(1-\gamma).
\]
The allocation is unchanged on the retained tuples. On any tuple the
sum of its full branch weights is at most one, since
$f_0+f_1=\rho\le1$ at every node of the choice tree. We may consequently
impose \eqref{ti:e-range} and reinstate all discarded tuples, with the
same negligible error \eqref{ti:large-group-error}.

\paragraph{Separation and its uniform cost.}
Use the Fourier convention
\[
 \widehat f(\xi)=\int_{\mathbb R}f(t)e^{-i\xi t}\,dt,
 \qquad f(t)=\int_{\mathbb R}\widehat f(\xi)e^{i\xi t}
                                    \frac{d\xi}{2\pi},
\]
and set
\[
 C_F=\max\left(1,
       \int|\widehat f_0(\xi)|\frac{d\xi}{2\pi},
       \int|\widehat f_1(\xi)|\frac{d\xi}{2\pi}\right).
\]
For a fixed branch $\delta=(\delta_1,\ldots,\delta_N)\in\{0,1\}^N$,
Fourier inversion of its $N$ transitions has total variation at most
$C_F^N$. Summing over the branches costs at most
\begin{equation}\label{ti:fourier-cost}
                         (2C_F)^N\le L^{C_f}
\end{equation}
for a fixed exponent, increased below to include the fixed cutoff
$\Psi$. In particular this exponent is independent of $K_0$.
Truncating each transition variable at $|\xi_i|\le L$ has total error
at most
\begin{equation}\label{ti:fourier-tail}
                  2^N N C_M C_F^{N-1}L^{-M}
\end{equation}
for any fixed $M$, by the union bound on the complementary integration
regions and the rapid decay of a single Fourier factor. The constant
$C_M$ occurs only once in each term of this estimate. Increasing $M$
therefore does not increase the exponent arising from $C_F^N$.

To see the separated phases explicitly, use
$R_i=\log E-\sum_{t<i}\delta_t\log p_t$. The product of the Fourier
phases is
\begin{equation}\label{ti:slot-frequencies}
 \exp\left(i\log E\sum_{i=1}^N\frac{\xi_i}{s_{j(i)}}\right)
       \prod_{t=1}^N p_t^{i\sigma_t},\qquad
 \sigma_t=-\frac{\xi_t}{s_{j(t)}}
           -\delta_t\sum_{i>t}\frac{\xi_i}{s_{j(i)}}.
\end{equation}
The dependence on $E$ is entirely in the scalar of modulus one.
Since $\sum_i s_{j(i)}^{-1}=O(T^{-1})$, the retained frequencies
satisfy $|\sigma_t|=O(L/T)$. Fourier inversion of the fixed smooth
function $t\mapsto\Psi(e^t)$ likewise separates $\Psi(eh/x)$,
with bounded integral norm; truncation at a fixed power of $L$ has
arbitrary logarithmic accuracy. Its contribution is a factor
$h^{i\varrho}$, a factor $e^{i\varrho}$ absorbed into the coefficient
of $e$, and a scalar of modulus one.

All these errors can be summed before Cauchy--Schwarz. The original
$m,n$ restrictions still imply $u\asymp x$, and the unmodified
factorization mass is bounded by
\[
 \sum_{\substack{mn=2u+1\\u\asymp x}}
       |\alpha_m\beta_n|A_0(u)
       \ll L^{O(1)}\sum_{v\asymp x}\tau(v)
       \ll xL^{O(1)}.
\]
Here and below $\tau$ is the divisor function. The same estimate
applies to the labelled-slot sum with its original factorial
normalization. Combining it with \eqref{ti:fourier-tail} proves an
error $O_H(xL^{-H})$ for arbitrary fixed $H$, without changing the
fixed exponent in \eqref{ti:fourier-cost}.

Consequently the required sum is, up to these errors, a sum and
integral of total variation at most $L^{C_f}$ of expressions
\begin{equation}\label{ti:separated-bilinear}
 \begin{split}
 B'&=\sum_{e,n}a_e(e)\beta_n
                  \sum_{mn=2eh+1}\alpha_m a_h(h),\\
 a_h(h)&=h^{i\varrho}\varepsilon(h)W_1(h)H(h_*),\qquad
 H(t)=\sum_{pr=t}\1_{p\in Q_{j_{**}}}p^{i\sigma}c_r.
 \end{split}
\end{equation}
The original $m,n$ restrictions are in $\alpha,\beta$, and the
$e$ range \eqref{ti:e-range} is in $a_e$. The prime in $H$ is the
designated prime. Its phase can be taken to be zero at this stage;
we allow a bounded-power frequency $\sigma$ in the notation. All
phases in \eqref{ti:separated-bilinear} have fixed logarithmic-power
bounds.

Both $a_e$ and $c_r$ are bounded coefficients. In fact, for a fixed
branch and a fixed product, disjointness of the bands determines the
prime multiset in every band. There are at most $r_0!$ assignments to
the relevant labelled slots of any one band, even when primes repeat.
Distributing the original factorial normalizations between the two
coefficients only improves the resulting bound
\begin{equation}\label{ti:coefficient-bounds}
 |a_e(e)|,|c_r|\le(r_0!)^{j_x+1}\le L^{C_a},
 \qquad |a_h(h)|\le L^{C_h}\tau(h_*).
\end{equation}
The last bound comes from summing over the designated prime divisor
of $h_*$. The exponents $C_a,C_h$ are fixed independently of $K_0$
and of the Fourier error accuracy. Only ordinary-band primes were
assigned to $e$, so its removal leaves the marked group weight and
the sign exactly on $h$, as used in \eqref{ti:separated-bilinear}.

\paragraph{Cauchy--Schwarz and the diagonal.}
Choose bounded nonnegative smooth majorants $\eta_e,\eta_n$, equal
to at least one on the $e,n$ ranges. We can arrange
\[
 \operatorname{supp}\eta_e\subset[\tfrac12L^{-C_s},2],
 \qquad \operatorname{supp}\eta_n\subset[1/2,3],
 \qquad
 \|\eta_e^{(j)}\|_\infty+\|\eta_n^{(j)}\|_\infty
                           \ll_j L^{C_{\eta}(j+1)},
\]
where $C_{\eta}$ depends only on $C_s$. There are $O(EV)$ pairs
$(e,n)$ in the original ranges. Hence \eqref{ti:coefficient-bounds}
and Cauchy--Schwarz give
\begin{equation}\label{ti:cauchy}
 |B'|^2\ll EVL^{C_p}\mathcal N,\qquad
 \mathcal N=\sum_{e,n}\eta_e(e/E)\eta_n(n/V)
                \left|\sum_{mn=2eh+1}\alpha_m a_h(h)\right|^2.
\end{equation}
Here $C_p$ is fixed independently of $K_0$.

The diagonal $m=m'$ in the square forces $h=h'$. Collecting by
$v=mn$ and then using \eqref{ti:coefficient-bounds} gives
\begin{align}
 \mathcal N_{\mathrm{diag}}
 &\ll L^{C_d}\sum_{\substack{v\asymp x\\v\ {\rm odd}}}
       \tau(v)\sum_{e\mid(v-1)/2}
                   \tau\left(\frac{v-1}{2e}\right)^2\notag\\
 &\le L^{C_d}\sum_{v\asymp x}\tau(v)\tau(v-1)^3
 \ll xL^{C_d'}.
 \label{ti:diagonal}
\end{align}
The final estimate follows from Cauchy--Schwarz and the fixed second
and sixth moments of $\tau$ in Proposition~\ref{pre:moments}.
Enlarging the $e$ divisor set in this
calculation makes $C_d'$ independent of $K_0$. Since
$EV\asymp xL^{-K_0}$, the diagonal contribution to the right side
of \eqref{ti:cauchy} is at most
$x^2L^{-K_0+C_p+C_d'}$. We now fix
\begin{equation}\label{ti:K-choice}
                  K_0>2(D_*+C_f+1)+C_p+C_d'.
\end{equation}
It remains to obtain an arbitrarily strong logarithmic saving for
the off-diagonal part of $\mathcal N$ with this fixed $K_0$.

\paragraph{The off-diagonal change of variables.}
A pair of representations in the expanded square satisfies
\[
                    mn=2eh+1,\qquad m'n=2eh'+1.
\]
The first identity gives $(n,2e)=1$. Subtracting the two identities
therefore gives a unique integer $k$ with
\begin{equation}\label{ti:determinant-forward}
 m'-m=2ek,\qquad h'-h=kn,\qquad
 z=m'h,\qquad z+k=mh'.
\end{equation}
Indeed $mh'-m'h=k(mn-2eh)=k$. On the off-diagonal $k\ne0$, and
the lower support bound for $e$ gives
\begin{equation}\label{ti:k-range}
                         0<|k|\ll L^{K_0+C_s}.
\end{equation}

For completeness, this parametrization is reversible for both signs
of $k$. Start from positive factorizations
$z=m'h$, $z+k=mh'$, require
$m'\equiv m\pmod{2|k|}$, and put
\begin{equation}\label{ti:determinant-reverse}
               e=\frac{m'-m}{2k}>0,\qquad
               n=\frac{2eh+1}{m}.
\end{equation}
The first quantity is an integer. The determinant identity implies
\[
                       m(h'-h)=k(2eh+1).
\]
Since $P^-(m)>W$ and $k$ has the fixed logarithmic-power bound
\eqref{ti:k-range}, we have $(m,k)=1$ for large $x$. Thus $n$ is
a positive integer and $h'-h=kn$. Substituting $m'=m+2ek$ recovers
$m'n=2eh'+1$. Every reconstructed variable is unique. Applying the
size cutoffs in \eqref{ti:cauchy} therefore gives an exact bijection,
with no additional multiplicity or divisibility condition.

Both $m,m'$ are odd and coprime to $k$. With $q_k=2|k|$, their
congruence is imposed exactly by
\begin{equation}\label{ti:character-orthogonality}
 \1_{m'\equiv m\ ({\rm mod}\ q_k)}
        =\frac1{\varphi(q_k)}\sum_{\chi\ ({\rm mod}\ q_k)}
                              \chi(m)\overline{\chi(m')}.
\end{equation}
The normalized character sum has total coefficient mass one.

\paragraph{The pulled-back cutoffs.}
On the supports above, $h\asymp x/e$ and $m'\asymp U$. Thus
$z\asymp Ux/e$ lies between fixed constant multiples of
$xL^{K_0}$ and $xL^{K_0+C_s}$. Insert a smooth dyadic partition
in $z$, retaining an outer cutoff $\psi_Z(z/Z)$ in every piece.
Only $O(1+T)$ dyads are needed, and all have $Z=xL^{O(1)}$.

For a fixed $k,Z$, put
\[
 t_1=\log(m/U),\quad t_2=\log(m'/U),\quad t_3=\log(z/Z).
\]
The arguments of the two majorants are exactly
\begin{equation}\label{ti:pulled-cutoffs}
 \begin{split}
 \frac eE
    &=\frac{m'-m}{2kE}
      =\frac{U}{2kE}(e^{t_2}-e^{t_1}),\\
 \frac nV
    &=\frac1V\left(\frac{(m'-m)z}{kmm'}+\frac1m\right)\\
    &=\frac{Z}{kUV}e^{t_3}(e^{-t_1}-e^{-t_2})
                                      +\frac1{UV}e^{-t_1}.
 \end{split}
\end{equation}
On a fixed compact box in these coordinates, every log-coordinate
derivative of the first expression is
$O_j(U/(|k|E))$, and every derivative of the second is
\[
                  O_j\left(\frac{Z}{|k|UV}+\frac1{UV}\right).
\]
All these quantities are bounded by fixed powers of $L$ after
\eqref{ti:K-choice}. The chain rule and the derivative bounds for
$\eta_e,\eta_n$ therefore give, for each multi-index $\nu$,
\begin{equation}\label{ti:pulled-derivatives}
             \|\partial^\nu b_{k,Z}\|_\infty
                         \ll_\nu L^{C_0(|\nu|+1)}
\end{equation}
for the pulled-back product of majorants, multiplied by fixed smooth
buffers in $t_1,t_2,t_3$. Choose these buffers to equal one on the
original $m,m'$ dyads and on the support of the outer $z$ cutoff.
This places $b_{k,Z}$ on a fixed compact box and leaves the sequence
$\alpha$ unchanged. The majorants vanish in neighborhoods of
nonpositive $e/E$ and $n/V$, so extending the product by zero across
those regions is smooth. Positivity in
\eqref{ti:determinant-reverse} is thereby imposed by the same cutoffs.

This is the fixed-dimensional setting of Lemma~\ref{pre:separation}.
Choose a fixed $C_8$ larger than the slope required in
\eqref{ti:pulled-derivatives}, and put $P=L^{C_8}$. Repeated
integration by parts in these three coordinates gives
\[
 |\widehat b_{k,Z}(\boldsymbol\xi)|
       \ll_j(1+|\boldsymbol\xi|/P)^{-j},\qquad
 \int_{\mathbb R^3}|\widehat b_{k,Z}(\boldsymbol\xi)|
                                      \,d\boldsymbol\xi\ll P^3,
\]
and
\begin{equation}\label{ti:post-fourier-tail}
 \int_{|\boldsymbol\xi|>PL}|\widehat b_{k,Z}(\boldsymbol\xi)|
                 \,d\boldsymbol\xi\ll_j P^3L^{3-j}.
\end{equation}
Consequently the cutoffs separate into phases
$m^{i\xi_1}(m')^{i\xi_2}z^{i\xi_3}$ at fixed logarithmic-power
cost and with fixed logarithmic-power frequencies. Scalars involving
$U$ and $Z$ have modulus one.

To justify the tails absolutely, an endpoint convolution is bounded
by $L^{O(1)}\tau(v_*)^2$: there are at most
$\tau_3(v_*)\le\tau(v_*)^2$ triples $mpr=v_*$. Including the bounded
group weights and dropping the congruence, Cauchy--Schwarz and the
fixed fourth divisor moment in Proposition~\ref{pre:moments} give
\begin{equation}\label{ti:absolute-correlation}
 L^{O(1)}\sum_{z\asymp Z}\tau(z)^2\tau(z+k)^2
                         \ll ZL^{O(1)}.
\end{equation}
Here $|k|\ll L^{O(1)}=o(Z)$ and both positions are positive.
This estimate controls \eqref{ti:post-fourier-tail}, after summing
over all $k$ and $Z$, to arbitrary logarithmic accuracy by increasing
$j$. The fixed exponent $C_8$ need not increase with that accuracy.

\paragraph{Identifying the two endpoints.}
Fix one character in \eqref{ti:character-orthogonality} and one
retained Fourier term. Before the character and Fourier factors, the
coefficient from the expanded square is
\[
              \alpha_m\overline{\alpha_{m'}}
                        a_h(h)\overline{a_h(h')}.
\]
Define invariant cores, for positive integers $v$, by
\begin{equation}\label{ti:endpoint-cores}
 \begin{split}
 F_0(v)&=\sum_{mpr=v_*}\alpha_m\chi(m)
             m^{i(\varrho-\xi_2)}
             \1_{p\in Q_{j_{**}}}p^{-i\sigma}\overline{c_r},\\
 G_0(v)&=\sum_{mpr=v_*}\alpha_m\chi(m)
             m^{i(\xi_1+\varrho)}
             \1_{p\in Q_{j_{**}}}p^{-i\sigma}\overline{c_r}.
 \end{split}
\end{equation}
Both cores contain the original restricted sequence $\alpha$.
All group primes are below $W$, whereas $m,m'$ have no prime factor
at most $W$. Therefore
\begin{equation}\label{ti:group-identities}
 \begin{gathered}
 z_*=m'h_*,\qquad (z+k)_*=mh'_*,\\
 W_1(z)=W_1(h),\quad W_1(z+k)=W_1(h'),\qquad
 \varepsilon(z)=\varepsilon(h),\quad
 \varepsilon(z+k)=\varepsilon(h').
 \end{gathered}
\end{equation}
Using also $h=z/m'$ and $h'=(z+k)/m$, direct expansion shows that
the separated sum is precisely
\begin{equation}\label{ti:separated-correlation}
 \sum_{\substack{z>0\\z+k>0}}
  \psi_Z(z/Z)z^{i(\varrho+\xi_3)}(z+k)^{-i\varrho}
       \overline{F_0(z)}G_0(z+k)
       W_1(z)W_1(z+k)\varepsilon(z)\varepsilon(z+k).
\end{equation}
For example, conjugating the first core supplies
$\overline{\alpha_{m'}}\overline{\chi(m')}
(m')^{i(\xi_2-\varrho)}H(h_*)$, while the second core supplies
$\alpha_m\chi(m)m^{i(\xi_1+\varrho)}\overline{H(h'_*)}$.
The remaining powers in \eqref{ti:separated-correlation} give exactly
$h^{i\varrho}(h')^{-i\varrho}z^{i\xi_3}$.
Thus the discrepancy-bearing first core has the required coefficient
$\alpha$, with its conjugation occurring outside that core.

Every hypothesis of Theorem~\ref{sh:correlation} is now satisfied.
The designated slot is a prime in the fixed range
\eqref{ti:designated-range}; the residual coefficients are bounded
by \eqref{ti:coefficient-bounds}; the character modulus $2|k|$ and
all frequencies have fixed logarithmic-power bounds; and the two
endpoints have the same allowed sign. Finally, set
\[
                    \widetilde\psi_Z(t)=t\psi_Z(t).
\]
Then $\psi_Z(z/Z)=(Z/z)\widetilde\psi_Z(z/Z)$. Consequently
\eqref{ti:separated-correlation} is exactly $Z$ times a harmonic
correlation covered by that theorem. For any prescribed $D_0$ it is
therefore
\begin{equation}\label{ti:correlation-bound}
                              O(ZL^{-D_0}),
\end{equation}
provided the discrepancy exponent $B$ is chosen sufficiently large
after all the fixed bounds and $D_0$.

\paragraph{Completing the parameter choices.}
Here is an explicit ordering of the choices. First fix the weight
geometry, $b_*,D_*,C$, the designated band, the allocation functions,
and the cutoff. These determine the exponents
\[
 C_s,\ C_a,\ C_h,\ C_f,\ C_p,\ C_d'.
\]
Fix $K_0$ by
\eqref{ti:K-choice}. This fixes the derivative and frequency bounds
after Cauchy--Schwarz, the Fourier scale $P$, and exponents
$K_k,K_z,C_{\rm sep}$ such that there are at most $L^{K_k}$ shifts,
$Z\le xL^{K_z}$, and the total post-Cauchy separation cost, including
the dyads, is at most $L^{C_{\rm sep}}$. These exponents are all fixed.

Choose $A_{\rm off}>0$ so that
\[
          -K_0+C_p-A_{\rm off}<-2(D_*+C_f+1),
\]
and then choose
\[
                 D_0>A_{\rm off}+K_k+K_z+C_{\rm sep}+2.
\]
Choose the differentiation orders in the Fourier tail estimates
large enough for the same error target.
Theorem~\ref{sh:correlation} now supplies a sufficiently large fixed $B$.
Summing \eqref{ti:correlation-bound} over shifts, dyads, characters,
and the retained Fourier integrals gives
\[
                       |\mathcal N_{\rm off}|
                                  \ll xL^{-A_{\rm off}}.
\]
Together with \eqref{ti:diagonal}, \eqref{ti:K-choice}, and
\eqref{ti:cauchy}, this proves
$|B'|\ll xL^{-D_*-C_f-1}$. The total pre-Cauchy variation
$L^{C_f}$ and the arbitrarily small previously discarded errors
prove \eqref{ti:bound} for each choice of $\varepsilon$, and hence
for $A$.
\end{proof}

\begin{remark}[Precision in subsequent proxy replacements]
\label{ti:proxy-order}
Theorem~\ref{ti:distribution} permits the actual interval restriction
to be included in $\alpha$, and it asks only for boundedness of the
other coefficient. Thus it can be applied a second time after
exchanging the two factors, even if the first coefficient has already
been replaced by a nonsmooth bounded proxy. On $O(L)$ contributing
factor dyads, its total error is $O(xL^{-D_*+1})$. The lower bound
$X_A\gg xL^{-C_A}$ in Lemma~\ref{wt:mass} makes this
$o(X_A/L)$ whenever $D_*>C_A+2$.

In particular, if a later proxy construction has coefficients bounded
by one independently of its logarithmic cell precision $S_1$, use
that coefficient bound when making all the choices above. First fix
$K_0$, the shifted-correlation accuracy, and finally the discrepancy
exponent $B$; only afterwards choose the proxy precision
$S_1>2B+3$. Neither $C_f$ nor the coefficient or diagonal exponents
depends on $S_1$. Increasing $S_1$ therefore changes only how large
$x$ must be and does not reopen any earlier choice.
\end{remark}
